\documentclass[12pt,a4paper]{article}

% ─── Codificación y lenguaje ─────────────────────────────────────────────────
\usepackage[utf8]{inputenc}
\usepackage[T1]{fontenc}
\usepackage[spanish,es-noquoting]{babel}

% ─── Diseño de página ────────────────────────────────────────────────────────
\usepackage[top=2.4cm, bottom=2.8cm, left=2.8cm, right=2.8cm]{geometry}
\usepackage{parskip}          % espaciado entre párrafos en lugar de sangría
\usepackage{microtype}        % mejora la justificación tipográfica
\usepackage{fancyhdr}
\pagestyle{fancy}
\fancyhf{}
\fancyhead[L]{\small\textit{IA Embebida en Sistemas Computacionales}}
\fancyhead[R]{\small\textit{Ingeniería Informática}}
\fancyfoot[C]{\thepage}
\renewcommand{\headrulewidth}{0.4pt}

% ─── Fuentes y tipografía ────────────────────────────────────────────────────
\usepackage{lmodern}
\usepackage{setspace}
\setstretch{1.15}

% ─── Colores ─────────────────────────────────────────────────────────────────
\usepackage{xcolor}
\definecolor{azulprimario}{RGB}{20, 70, 140}
\definecolor{azulclaro}{RGB}{230, 240, 255}
\definecolor{verdeok}{RGB}{10, 110, 50}
\definecolor{verdeclaro}{RGB}{230, 248, 236}
\definecolor{rojowarn}{RGB}{160, 30, 20}
\definecolor{rojoclaro}{RGB}{255, 235, 230}
\definecolor{naranjainf}{RGB}{160, 90, 0}
\definecolor{naranjaclaro}{RGB}{255, 245, 220}
\definecolor{grisback}{RGB}{248, 248, 246}
\definecolor{codegreen}{rgb}{0,0.52,0}
\definecolor{codegray}{rgb}{0.5,0.5,0.5}
\definecolor{codepurple}{rgb}{0.5,0,0.62}
\definecolor{codeblue}{rgb}{0.0,0.18,0.72}
\definecolor{backcolour}{rgb}{0.96,0.96,0.94}

% ─── Cajas de contenido (tcolorbox) ──────────────────────────────────────────
\usepackage{tcolorbox}
\tcbuselibrary{skins, breakable}

\newtcolorbox{conceptbox}[1]{
  title={\textbf{#1}},
  colback=azulclaro, colframe=azulprimario,
  fonttitle=\color{white}, coltitle=white,
  colbacktitle=azulprimario,
  attach boxed title to top left={yshift=-2mm, xshift=4mm},
  boxed title style={rounded corners, size=small},
  rounded corners, breakable, left=6pt, right=6pt,
  top=8pt, bottom=6pt,
}

\newtcolorbox{warnbox}[1]{
  title={\textbf{#1}},
  colback=rojoclaro, colframe=rojowarn,
  fonttitle=\color{white}, coltitle=white,
  colbacktitle=rojowarn,
  attach boxed title to top left={yshift=-2mm, xshift=4mm},
  boxed title style={rounded corners, size=small},
  rounded corners, breakable, left=6pt, right=6pt,
  top=8pt, bottom=6pt,
}

\newtcolorbox{examplebox}[1]{
  title={\textbf{#1}},
  colback=verdeclaro, colframe=verdeok,
  fonttitle=\color{white}, coltitle=white,
  colbacktitle=verdeok,
  attach boxed title to top left={yshift=-2mm, xshift=4mm},
  boxed title style={rounded corners, size=small},
  rounded corners, breakable, left=6pt, right=6pt,
  top=8pt, bottom=6pt,
}

\newtcolorbox{infobox}[1]{
  title={\textbf{#1}},
  colback=naranjaclaro, colframe=naranjainf,
  fonttitle=\color{white}, coltitle=white,
  colbacktitle=naranjainf,
  attach boxed title to top left={yshift=-2mm, xshift=4mm},
  boxed title style={rounded corners, size=small},
  rounded corners, breakable, left=6pt, right=6pt,
  top=8pt, bottom=6pt,
}

% ─── Código fuente ───────────────────────────────────────────────────────────
\usepackage{listings}

\lstdefinestyle{pythonstyle}{
  backgroundcolor=\color{backcolour},
  commentstyle=\color{codegreen}\itshape,
  keywordstyle=\color{codeblue}\bfseries,
  numberstyle=\tiny\color{codegray},
  stringstyle=\color{codepurple},
  basicstyle=\ttfamily\small,
  breakatwhitespace=false, breaklines=true,
  keepspaces=true, numbers=left, numbersep=8pt,
  showstringspaces=false, tabsize=2,
  language=Python, frame=single,
  rulecolor=\color{black!25}, xleftmargin=14pt,
  morekeywords={from,import,as,def,return,True,False,
                None,class,self,with,yield,lambda,async,await},
  captionpos=b,
}

\lstdefinestyle{jsonstyle}{
  backgroundcolor=\color{backcolour},
  basicstyle=\ttfamily\small,
  breaklines=true, frame=single,
  rulecolor=\color{black!25},
  showstringspaces=false, xleftmargin=14pt,
  morestring=[b]", stringstyle=\color{codepurple},
  captionpos=b,
}

\lstdefinestyle{shellstyle}{
  backgroundcolor=\color{gray!10},
  basicstyle=\ttfamily\small,
  breaklines=true, frame=single,
  rulecolor=\color{black!25},
  showstringspaces=false, xleftmargin=14pt,
}

% ─── Diagramas TikZ ──────────────────────────────────────────────────────────
\usepackage{tikz}
\usetikzlibrary{shapes.geometric, arrows.meta, positioning,
                fit, backgrounds, calc, shadows}

% ─── Otros paquetes ──────────────────────────────────────────────────────────
\usepackage{booktabs}
\usepackage{amsmath}
\usepackage{enumitem}
\usepackage{float}
\usepackage{caption}
\usepackage{subcaption}
\usepackage{graphicx}
\usepackage[hidelinks, colorlinks=true,
            linkcolor=azulprimario,
            urlcolor=azulprimario]{hyperref}

% ─── Títulos de sección ───────────────────────────────────────────────────────
\usepackage{titlesec}
\titleformat{\section}
  {\Large\bfseries\color{azulprimario}}{\thesection.}{0.6em}{}[\vspace{-0.4em}\rule{\linewidth}{0.5pt}]
\titleformat{\subsection}
  {\large\bfseries\color{azulprimario!80!black}}{\thesubsection.}{0.5em}{}
\titleformat{\subsubsection}
  {\normalsize\bfseries}{\thesubsubsection.}{0.4em}{}

% ─── Información del documento ───────────────────────────────────────────────
\title{
  \vspace{-1cm}
  {\large Ingeniería Informática --- Material de Clase}\\[0.5em]
  {\LARGE\bfseries\color{azulprimario}
    IA Embebida en Sistemas Computacionales}\\[0.4em]
  {\large LLM \textbullet{} RAG \textbullet{} Redis \textbullet{}
   LangChain \textbullet{} Ollama}
}
\author{Departamento de Ciencias de la Computación}
\date{Semestre I --- 2026}

% ─────────────────────────────────────────────────────────────────────────────
\begin{document}

\maketitle
\thispagestyle{fancy}

\begin{abstract}
Este documento desarrolla los conceptos y la implementación práctica de un sistema
de inteligencia artificial embebido en una arquitectura de software real. El caso
de estudio es un sistema de generación automática de minutas alimentarias para
niños menores de 3 años con alergias alimentarias. Se estudia la integración de
Ollama como motor de inferencia local, LangChain como orquestador de componentes,
Redis como memoria contextual de sesión y RAG como mecanismo de recuperación
documental. El objetivo central es comprender el LLM como un \textbf{componente}
dentro de un sistema mayor, no como la aplicación en sí misma.

\medskip
\noindent
\textbf{Prerrequisitos:} Modelos LLM locales con Ollama, embeddings, chunking,
RAG básico y recuperación documental.
\end{abstract}

\tableofcontents
\vspace{1em}
\rule{\linewidth}{0.5pt}

% ═══════════════════════════════════════════════════════════════════
\section{Motivación: Del LLM Aislado al Sistema Inteligente}
% ═══════════════════════════════════════════════════════════════════

El error conceptual más frecuente en proyectos de IA aplicada es tratar al
LLM como si fuera la aplicación completa. Un modelo de lenguaje no es una
solución por sí mismo: es un componente que necesita memoria, contexto de
dominio, restricciones estructurales e integración con otros sistemas para
producir valor real en producción.

Esta clase propone una narrativa de evolución arquitectural que iremos
construyendo paso a paso:

\begin{center}
\large
\textbf{LLM solo}
$\;\longrightarrow\;$
\textbf{LLM + RAG}
$\;\longrightarrow\;$
\textbf{Sistema inteligente}
$\;\longrightarrow\;$
\textbf{Aplicación empresarial}
\end{center}

\begin{warnbox}{El problema de tratar el LLM como solución completa}
Un LLM llamado directamente sin arquitectura:
\begin{itemize}[noitemsep]
  \item No recuerda conversaciones anteriores (stateless por defecto)
  \item Desconoce la normativa específica del dominio
  \item Produce respuestas en texto libre, imposibles de parsear automáticamente
  \item Alucina valores, cantidades e ingredientes inexistentes
  \item Es imposible de auditar o rastrear
\end{itemize}
\end{warnbox}

\begin{conceptbox}{Principio de diseño central}
El LLM es uno de los componentes del sistema, con responsabilidades
bien definidas. LangChain lo orquesta, RAG le provee conocimiento del
dominio, Redis le da memoria de sesión, y el parser garantiza la
estructura de salida.
\end{conceptbox}

% ═══════════════════════════════════════════════════════════════════
\section{Caso de Estudio: Sistema de Minutas Alimentarias}
% ═══════════════════════════════════════════════════════════════════

\subsection{Contexto del problema}

Los jardines infantiles JUNJI y VTF atienden a niños de 0 a 3 años cuyos
perfiles nutricionales están documentados, incluyendo alergias e intolerancias
específicas. Actualmente, una nutricionista dedica aproximadamente 8 horas
semanales a generar minutas personalizadas para cada niño, un proceso
repetitivo y propenso a errores humanos con consecuencias médicas graves.

El sistema que construiremos automatiza la generación de estas minutas
respetando las siguientes restricciones del dominio:

\begin{itemize}
  \item \textbf{Normativa JUNJI circular N$^{\circ}$73 (2022):} porciones
    por grupo alimentario y rangos calóricos por tiempo de comida.
  \item \textbf{Tablas MINSAL de composición de alimentos:} valores
    nutricionales por ingrediente, con nombres en español chileno.
  \item \textbf{Restricción calórica:} 200--350 kcal por tiempo de comida
    para niños de 12 a 36 meses.
  \item \textbf{Alergias declaradas:} ningún ingrediente del plato puede
    contener el alérgeno declarado, ni siquiera como trazas.
\end{itemize}

\subsection{Salida esperada del sistema}

El sistema no debe retornar texto libre. Debe retornar un objeto JSON
estructurado, validado y directamente consumible por otras aplicaciones
(app móvil de padres, sistema de compras ERP, tickets de cocina, calendario
institucional). Esta exigencia de estructura es lo que diferencia un prototipo
de demostración de un sistema en producción.

La figura \ref{fig:json_esquema} muestra el esquema de la respuesta esperada,
que se detalla en la sección \ref{sec:json}.

% ═══════════════════════════════════════════════════════════════════
\section{Arquitecturas: Tradicional vs. Moderna}
\label{sec:arquitecturas}
% ═══════════════════════════════════════════════════════════════════

\subsection{Arquitectura tradicional: el enfoque incorrecto}

El enfoque más intuitivo --- y el más problemático --- es llamar al LLM
directamente desde la aplicación con un prompt construido en el momento:

\begin{figure}[H]
\centering
\begin{tikzpicture}[
  box/.style={rectangle, rounded corners=4pt, minimum width=2.2cm,
              minimum height=0.8cm, text centered, font=\small,
              draw=black!50, fill=white, drop shadow},
  boxred/.style={rectangle, rounded corners=4pt, minimum width=2.2cm,
              minimum height=0.8cm, text centered, font=\small,
              draw=rojowarn, fill=rojoclaro!60, drop shadow},
  arr/.style={-Stealth, thick, draw=black!55},
  arrd/.style={-Stealth, thick, draw=rojowarn, dashed},
]
  \node[box] (user) at (0,0)   {Usuario};
  \node[box] (app)  at (3.2,0)   {Aplicación};
  \node[boxred] (llm) at (6.4,0)   {LLM (Ollama)};
  \node[boxred] (resp) at (9.6,0) {Texto libre};

  \draw[arr] (user) -- (app);
  \draw[arr] (app)  -- node[above, font=\scriptsize]{prompt directo} (llm);
  \draw[arr] (llm)  -- (resp);

  \node[boxred, text width=2.6cm, align=center, font=\scriptsize] (p1) at (2.0,-2.0)
        {\textbf{Sin memoria}\\ sesión stateless};
  \node[boxred, text width=2.6cm, align=center, font=\scriptsize] (p2) at (5.5,-2.0)
        {\textbf{Sin contexto}\\ ignora normativa};
  \node[boxred, text width=2.6cm, align=center, font=\scriptsize] (p3) at (9.0,-2.0)
        {\textbf{Sin estructura}\\ texto no parseable};

  \draw[arrd] (app.south)  -- (p1.north);
  \draw[arrd] (llm.south)  -- (p2.north);
  \draw[arrd] (resp.south) -- (p3.north);
\end{tikzpicture}
\caption{Arquitectura tradicional y sus tres puntos de fallo}
\label{fig:arq_tradicional}
\end{figure}

Esta arquitectura falla en tres dimensiones independientes:

\begin{description}
  \item[Problema 1 --- Memoria.] Cada invocación es un contexto nuevo
    e independiente. El LLM no recuerda las alergias de María entre
    la consulta del lunes y la del martes. El nutricionista debe
    repetir el perfil completo en cada sesión. A 50 niños $\times$
    3 tiempos $\times$ 5 días, esto es inmanejable.

  \item[Problema 2 --- Contexto de dominio.] Ningún LLM de propósito
    general fue entrenado con la circular JUNJI N$^{\circ}$73 de 2022
    ni con las tablas de composición MINSAL vigentes. Las porciones
    generadas son genéricas, los valores calóricos son aproximados y
    los ingredientes pueden ser nombres extranjeros inusuales en Chile.

  \item[Problema 3 --- Integración.] La respuesta en texto libre tiene
    un formato inconsistente entre ejecuciones. Es imposible conectar
    esta salida automáticamente con un ERP, una app móvil o un sistema
    de tickets de cocina sin un frágil parser basado en expresiones
    regulares que se rompe ante cualquier variación del modelo.
\end{description}

\begin{warnbox}{Conclusión del análisis}
Necesitamos una arquitectura que provea simultáneamente:
\textbf{memoria} (Redis) + \textbf{contexto documental} (RAG) +
\textbf{salida estructurada} (JSON + Pydantic).
\end{warnbox}

\subsection{Arquitectura moderna basada en IA}

La arquitectura moderna separa responsabilidades de manera explícita.
Cada componente resuelve exactamente un problema:

\begin{figure}[H]
\centering
\begin{tikzpicture}[
  box/.style={rectangle, rounded corners=4pt, minimum width=2.0cm,
              minimum height=0.7cm, text centered, font=\small,
              draw=azulprimario!60, fill=azulclaro, drop shadow},
  boxg/.style={rectangle, rounded corners=4pt, minimum width=2.0cm,
              minimum height=0.7cm, text centered, font=\small,
              draw=verdeok, fill=verdeclaro, drop shadow},
  boxo/.style={rectangle, rounded corners=4pt, minimum width=2.0cm,
              minimum height=0.7cm, text centered, font=\small,
              draw=naranjainf, fill=naranjaclaro, drop shadow},
  boxp/.style={rectangle, rounded corners=4pt, minimum width=2.0cm,
              minimum height=0.7cm, text centered, font=\small,
              draw=purple!65, fill=purple!8, drop shadow},
  arr/.style={-Stealth, thick, draw=gray!65},
  arrd/.style={-Stealth, thick, draw=purple!55, dashed},
  lc/.style={rectangle, rounded corners=6pt, draw=azulprimario!40,
             fill=azulclaro!40, dashed, minimum width=8.0cm,
             minimum height=1.6cm},
]
  \node[lc] at (6.2, 0) {};
  \node[above, font=\small\bfseries\color{azulprimario}] at (6.2, 0.8)
    {LangChain LCEL Chain};

  \node[box]  (prompt)  at (3.5, 0) {Prompt\\Template};
  \node[box]  (ollama)  at (6.2, 0) {Ollama\\LLaMA 3.2};
  \node[boxg] (parser)  at (9.0, 0) {JSON\\Parser};

  \node[boxo] (user)    at (0.5, 0) {FastAPI};
  \node[boxp] (rag)     at (3.5,-1.8) {RAG (Chroma)};
  \node[boxp] (redis)   at (6.2,-1.8) {Redis (sesión)};
  \node[boxg] (output)  at (12.0, 0) {JSON\\Minuta};

  \draw[arr] (user)   -- (prompt);
  \draw[arr] (prompt) -- (ollama);
  \draw[arr] (ollama) -- (parser);
  \draw[arr, very thick, draw=verdeok] (parser) -- (output);
  \draw[arrd] (rag)   -- node[right, font=\scriptsize]{contexto normativo} (prompt);
  \draw[arrd] (redis) -- node[right, font=\scriptsize]{historial sesión} (ollama);
\end{tikzpicture}
\caption{Arquitectura moderna: LLM como componente orquestado}
\label{fig:arq_moderna}
\end{figure}

La tabla \ref{tab:componentes} resume las responsabilidades de cada
componente en el sistema.

\begin{table}[H]
\centering
\caption{Responsabilidades de cada componente}
\label{tab:componentes}
\begin{tabular}{lll}
\toprule
\textbf{Componente} & \textbf{Problema que resuelve} & \textbf{Tecnología} \\
\midrule
Ollama + LLaMA 3.2 & Inferencia de lenguaje & Motor llama.cpp local \\
LangChain (LCEL)   & Orquestación y composición & Runnable pipeline \\
RAG + Chroma       & Contexto documental del dominio & Vector store + MMR \\
Redis              & Memoria de sesión persistente & In-memory store \\
JSON Parser        & Estructura de salida garantizada & grammar-guided sampling \\
Pydantic v2        & Validación y tipado de la salida & Schema enforcement \\
FastAPI            & Exposición como API REST & Endpoint HTTP \\
\bottomrule
\end{tabular}
\end{table}

% ═══════════════════════════════════════════════════════════════════
\section{Componente RAG: Retrieval-Augmented Generation}
\label{sec:rag}
% ═══════════════════════════════════════════════════════════════════

\subsection{Definición y motivación}

\begin{conceptbox}{RAG -- Definición}
RAG (\textit{Retrieval-Augmented Generation}) es un patrón arquitectural
que \textbf{aumenta dinámicamente} el contexto del prompt con fragmentos
recuperados de una base de conocimiento externa, sin modificar los pesos
del modelo. El LLM razona sobre información que no conocía en entrenamiento.
\end{conceptbox}

En nuestro caso, el corpus normativo (circular JUNJI, tablas MINSAL) cambia
periódicamente. Con RAG, basta con re-indexar los documentos actualizados sin
reentrenar el modelo. Esta capacidad de actualización sin reentrenamiento es
la ventaja más importante frente al fine-tuning.

\begin{table}[H]
\centering
\caption{Comparación: RAG vs. Fine-tuning}
\begin{tabular}{lcc}
\toprule
\textbf{Criterio} & \textbf{RAG} & \textbf{Fine-tuning} \\
\midrule
Actualizar conocimiento & Solo re-indexar & Reentrenar ($\$$\$) \\
Trazabilidad de respuestas & Sí (se conoce el chunk) & No \\
GPU para entrenamiento & No & Sí \\
Costo operativo & Bajo & Alto \\
Auditable legalmente & Sí & Difícil \\
Latencia adicional & $\sim$50--150 ms & 0 ms \\
\bottomrule
\end{tabular}
\end{table}

\begin{infobox}{Limitación importante de RAG}
RAG \textbf{no mejora el razonamiento} del LLM ni su capacidad de seguir
instrucciones. Únicamente amplía la base de conocimiento accesible en
tiempo de inferencia. Si el LLM base razona mal, RAG no lo corrige.
\end{infobox}

\subsection{Pipeline de RAG: dos fases}

El pipeline de RAG se divide en dos fases con ciclos de vida distintos:

\begin{figure}[H]
\centering
\begin{tikzpicture}[
  doc/.style={rectangle, rounded corners=2pt, minimum width=1.7cm,
              minimum height=0.6cm, text centered, font=\small,
              draw=naranjainf, fill=naranjaclaro},
  proc/.style={rectangle, rounded corners=2pt, minimum width=1.9cm,
               minimum height=0.6cm, text centered, font=\small,
               draw=azulprimario!60, fill=azulclaro},
  db/.style={cylinder, shape border rotate=90, aspect=0.35,
             minimum width=1.7cm, minimum height=0.75cm,
             text centered, font=\small,
             draw=purple!60, fill=purple!8},
  arr/.style={-Stealth, thick, draw=gray!65},
  arrd/.style={-Stealth, dashed, draw=purple!55, thick},
  lbl/.style={font=\small\bfseries, text=gray!60},
]
  \node[lbl] at (0.8, 2.1) {FASE 1 --- Indexación (offline, ejecutar una vez)};
  \node[doc]  (pdf)    at (0.0, 1.35) {PDFs\\Normativa};
  \node[proc] (chunk)  at (2.4, 1.35) {Chunking\\512 tokens};
  \node[proc] (embed)  at (4.9, 1.35) {Encoder\\nomic-embed};
  \node[db]   (chroma) at (7.4, 1.35) {Chroma DB\\vectores};

  \draw[arr] (pdf)   -- (chunk);
  \draw[arr] (chunk) -- (embed);
  \draw[arr] (embed) -- node[above, font=\scriptsize]{float32 768-dim} (chroma);

  \node[lbl] at (0.8, 0.55) {FASE 2 --- Recuperación (online, cada consulta)};
  \node[doc]  (query) at (0.0,-0.2) {Consulta};
  \node[proc] (qemb)  at (2.4,-0.2) {Encoder\\(mismo)};
  \node[proc] (sim)   at (4.9,-0.2) {Similitud\\coseno (MMR)};
  \node[proc] (topk)  at (7.4,-0.2) {Top-4\\docs};
  \node[proc] (augp)  at (9.9,-0.2) {Prompt\\aumentado};
  \node[doc]  (out)   at (12.4,-0.2){LLM $\to$\\JSON};

  \draw[arr] (query) -- (qemb);
  \draw[arr] (qemb)  -- (sim);
  \draw[arr] (sim)   -- (topk);
  \draw[arr] (topk)  -- (augp);
  \draw[arr, very thick, draw=azulprimario] (augp) -- (out);
  \draw[arrd] (chroma.south) -- node[right, font=\scriptsize]{ANN search} (sim.north);
\end{tikzpicture}
\caption{Pipeline completo de RAG: fases de indexación y recuperación}
\label{fig:rag_pipeline}
\end{figure}

\begin{conceptbox}{Invariante del embedding}
El \textbf{mismo modelo de embedding} debe usarse en la fase de indexación
y en la fase de recuperación. Mezclar modelos hace que los vectores queden
en espacios semánticos incompatibles y el retriever retorna basura.
Para este sistema usamos \texttt{nomic-embed-text} (768 dimensiones) via Ollama.
\end{conceptbox}

\subsection{Implementación del indexador RAG}

\begin{lstlisting}[style=pythonstyle, caption={rag\_indexer.py --- Indexación del corpus normativo}]
from langchain_community.document_loaders import DirectoryLoader, PyPDFLoader
from langchain.text_splitter import RecursiveCharacterTextSplitter
from langchain_chroma import Chroma
from config import embeddings, CHROMA_DIR

def build_vector_store(pdf_dir: str) -> Chroma:
    """Indexa el corpus normativo completo en Chroma."""
    loader = DirectoryLoader(
        pdf_dir,
        glob="**/*.pdf",
        loader_cls=PyPDFLoader,
    )
    docs = loader.load()
    print(f"Documentos cargados: {len(docs)}")

    splitter = RecursiveCharacterTextSplitter(
        chunk_size=512,
        chunk_overlap=64,
        separators=["\n\n", "\n", ". ", " ", ""],
        length_function=len,
    )
    chunks = splitter.split_documents(docs)
    print(f"Fragmentos generados: {len(chunks)}")

    vectorstore = Chroma.from_documents(
        documents=chunks,
        embedding=embeddings,
        persist_directory=CHROMA_DIR,
        collection_name="normativa_junji",
    )
    return vectorstore

def get_retriever(vectorstore: Chroma):
    """Retriever con Max Marginal Relevance para diversidad de fragmentos."""
    return vectorstore.as_retriever(
        search_type="mmr",
        search_kwargs={
            "k": 4,           # fragmentos a retornar
            "fetch_k": 20,    # candidatos iniciales
            "lambda_mult": 0.7,  # balance relevancia/diversidad
        },
    )

if __name__ == "__main__":
    vs = build_vector_store("./documentos_normativos")
    print("Indexacion completada.")
\end{lstlisting}

\begin{infobox}{MMR vs. Similarity Search}
\textbf{Similarity Search} retorna los 4 fragmentos más similares a la
consulta, que pueden ser 4 oraciones del mismo párrafo del mismo documento.
\textbf{MMR (Max Marginal Relevance)} balancea relevancia y \textit{diversidad}:
evita recuperar fragmentos redundantes. Para corpus normativo donde la misma
norma puede estar parafraseada varias veces, MMR produce un contexto más rico.
\end{infobox}

% ═══════════════════════════════════════════════════════════════════
\section{Componente Redis: Memoria Contextual de Sesión}
\label{sec:redis}
% ═══════════════════════════════════════════════════════════════════

\subsection{Por qué Redis y no una base de datos relacional}

Un LLM necesita que el historial de conversación sea inyectado en el
prompt en cada invocación. Esto implica leer el historial con muy baja
latencia antes de cada inferencia. Redis es la elección natural por
varias razones:

\begin{itemize}
  \item \textbf{In-memory}: latencia $<$ 1 ms, frente a 5--20 ms de
    PostgreSQL para lecturas de registros.
  \item \textbf{TTL nativo}: las sesiones expiran automáticamente sin
    necesidad de cron jobs de limpieza.
  \item \textbf{Tipo lista}: el historial de mensajes se almacena como
    una lista Redis, permitiendo operaciones de recorte (\texttt{LTRIM})
    para implementar ventana deslizante.
  \item \textbf{Integración directa}: LangChain provee
    \texttt{RedisChatMessageHistory} que serializa/deserializa el
    historial automáticamente.
\end{itemize}

\subsection{Tipos de memoria disponibles en LangChain}

\begin{table}[H]
\centering
\caption{Estrategias de memoria para LLMs en LangChain}
\begin{tabular}{lll}
\toprule
\textbf{Clase} & \textbf{Estrategia} & \textbf{Uso recomendado} \\
\midrule
\texttt{ConversationBufferMemory} & Historia completa & Sesiones cortas ($<$20 turns) \\
\texttt{ConversationWindowMemory} & Últimas $N$ interacciones & Sesiones largas con presupuesto fijo \\
\texttt{ConversationSummaryMemory} & Resumen comprimido por LLM & Conversaciones muy largas \\
\texttt{ConversationEntityMemory} & Entidades nombradas & Seguimiento de personajes/datos \\
\texttt{VectorStoreRetrieverMemory} & Búsqueda semántica en el historial & Memorias muy largas \\
\bottomrule
\end{tabular}
\end{table}

\begin{warnbox}{Trade-off crítico: tokens vs. contexto}
Más historial en Redis equivale a mejor contexto para el LLM, pero también
a más tokens en el prompt en cada invocación, mayor latencia de inferencia
y mayor consumo de VRAM. Para \texttt{llama3.2:3b} con ventana de 4096 tokens,
reservar máximo 1500 tokens para historial, 1000 para contexto RAG y
1500 para la respuesta.
\end{warnbox}

\subsection{Implementación de la memoria con Redis}

\begin{lstlisting}[style=pythonstyle, caption={memory.py --- Memoria de sesión con Redis}]
from langchain_community.chat_message_histories import RedisChatMessageHistory
from langchain_core.runnables.history import RunnableWithMessageHistory

REDIS_URL = "redis://localhost:6379"

def get_session_history(session_id: str) -> RedisChatMessageHistory:
    """
    Factory de historiales: cada session_id tiene su propia clave en Redis.
    Esquema de clave recomendado: "jardin:{id}:nino:{id}:fecha:{YYYY-MM-DD}"
    """
    return RedisChatMessageHistory(
        session_id=session_id,
        url=REDIS_URL,
        ttl=86400,          # 24 horas en segundos
        key_prefix="minuta:",
    )

# Envolver la cadena base con memoria persistente
def add_memory_to_chain(base_chain):
    return RunnableWithMessageHistory(
        base_chain,
        get_session_history,
        input_messages_key="input",
        history_messages_key="history",
    )
\end{lstlisting}

La clave de sesión en Redis sigue este esquema:
\begin{center}
\texttt{minuta:jardin01:nino:maria\_001}
\end{center}

Cada entrada en la lista Redis tiene la siguiente estructura JSON:

\begin{lstlisting}[style=jsonstyle, caption={Estructura de un mensaje almacenado en Redis}]
{
  "type": "human",
  "content": "Maria, 2 anios, alergia a lactosa. Genera almuerzo del lunes.",
  "additional_kwargs": {},
  "response_metadata": {}
}
\end{lstlisting}

En la segunda invocación para María, LangChain lee esta lista automáticamente
y la inyecta en el placeholder \texttt{\{history\}} del prompt. El LLM ya
sabe que María es alérgica a la lactosa \textbf{sin que el usuario lo repita}.

% ═══════════════════════════════════════════════════════════════════
\section{Componente LangChain: Orquestación con LCEL}
\label{sec:langchain}
% ═══════════════════════════════════════════════════════════════════

\subsection{LangChain Expression Language (LCEL)}

LCEL define las cadenas de procesamiento como \textbf{composición funcional}
de objetos \texttt{Runnable} mediante el operador \texttt{|}. Cada
\texttt{Runnable} expone la misma interfaz: \texttt{invoke()},
\texttt{stream()}, \texttt{batch()} y sus variantes asíncronas.

\begin{equation}
\text{chain} = \underbrace{\text{prompt}}_{\text{Runnable}} \;|\;
               \underbrace{\text{llm}}_{\text{Runnable}} \;|\;
               \underbrace{\text{parser}}_{\text{Runnable}}
\end{equation}

Esta composición es análoga a una pipeline Unix:
\texttt{cat archivo | grep patrón | sort}. Cada componente recibe
la salida del anterior y produce su propia salida.

\begin{conceptbox}{Ventajas de LCEL sobre llamadas directas a Ollama}
\begin{itemize}[noitemsep]
  \item \textbf{Retry automático:} reintenta ante fallos de red o timeout
  \item \textbf{Streaming:} tokens disponibles en tiempo real sin cambiar el código
  \item \textbf{Observabilidad:} integración nativa con LangSmith para tracing
  \item \textbf{Portabilidad:} cambiar de \texttt{ChatOllama} a \texttt{ChatOpenAI}
    requiere modificar una sola línea, no toda la lógica
  \item \textbf{Paralelismo:} \texttt{RunnableParallel} ejecuta ramas concurrentemente
\end{itemize}
\end{conceptbox}

\begin{warnbox}{Versión importante}
Esta clase usa \texttt{langchain >= 0.3.x} con LCEL nativo.
La API anterior basada en \texttt{LLMChain} está \textbf{deprecada}.
No usar ejemplos de tutoriales previos a 2024.
\end{warnbox}

\subsection{Construcción de la cadena LCEL}

\begin{lstlisting}[style=pythonstyle, caption={chain\_builder.py --- Cadena LCEL con RAG, memoria y parser JSON}]
from langchain_ollama import ChatOllama
from langchain_core.prompts import ChatPromptTemplate
from langchain_core.output_parsers import JsonOutputParser
from langchain_core.runnables import RunnablePassthrough
from langchain_core.runnables.history import RunnableWithMessageHistory
from operator import itemgetter
from memory import get_session_history

# ── Configuracion del LLM ───────────────────────────────────────────────────
llm = ChatOllama(
    model="llama3.2",
    temperature=0.1,   # baja temperatura = respuestas mas consistentes
    format="json",     # grammar-guided sampling en llama.cpp
    num_ctx=4096,
)

# ── Plantilla del prompt ────────────────────────────────────────────────────
prompt = ChatPromptTemplate.from_messages([
    ("system",      "{system_prompt}"),
    ("placeholder", "{history}"),       # <- inyectado por Redis
    ("human",
     "Contexto normativo recuperado:\n{context}\n\n"
     "Solicitud: {input}"),
])

# ── Cadena base: RAG + Prompt + LLM + Parser ───────────────────────────────
base_chain = (
    RunnablePassthrough.assign(
        context=(itemgetter("input") | retriever)
    )
    | prompt
    | llm
    | JsonOutputParser()
)

# ── Cadena final: con memoria Redis ────────────────────────────────────────
chain = RunnableWithMessageHistory(
    base_chain,
    get_session_history,
    input_messages_key="input",
    history_messages_key="history",
)

# ── Invocacion ──────────────────────────────────────────────────────────────
result: dict = chain.invoke(
    {
        "input": "Genera almuerzo del lunes para Maria, 2 anios, alergia lactosa",
        "system_prompt": SYSTEM_PROMPT,
    },
    config={"configurable": {"session_id": "jard01:maria"}},
)
# result es un dict Python validado, listo para JSONResponse en FastAPI
\end{lstlisting}

\subsubsection{Flujo de datos paso a paso}

\begin{enumerate}
  \item \textbf{Input}: el texto de la consulta del nutricionista llega
    como string.
  \item \textbf{RunnablePassthrough.assign(context=...)}: el retriever
    de Chroma es invocado con la consulta y retorna los 4 fragmentos
    normativos más relevantes. Estos se asignan al campo \texttt{context}
    del diccionario.
  \item \textbf{ChatPromptTemplate}: ensambla el mensaje final combinando
    el system prompt, el historial de Redis, el contexto normativo y la
    consulta del usuario.
  \item \textbf{ChatOllama}: Ollama ejecuta la inferencia local con
    grammar-guided sampling, garantizando que los tokens generados
    formen JSON sintácticamente válido.
  \item \textbf{JsonOutputParser}: convierte el string JSON a un
    \texttt{dict} Python, que es el tipo de retorno de \texttt{invoke()}.
\end{enumerate}

% ═══════════════════════════════════════════════════════════════════
\section{Pre-Prompt Engineering}
\label{sec:prompt}
% ═══════════════════════════════════════════════════════════════════

\subsection{Qué es y por qué importa}

El \textit{system prompt} define el rol, las restricciones y el comportamiento
del LLM dentro del sistema. Es la interfaz de programación del modelo: no
es una sugerencia, sino un conjunto de \textbf{reglas de ejecución} que el
modelo debe seguir.

Un system prompt bien diseñado para un sistema de producción debe incluir:
\begin{enumerate}
  \item Definición del rol y las credenciales de dominio
  \item Reglas absolutas numeradas y sin ambigüedad
  \item Parámetros cuantitativos precisos (rangos calóricos, unidades)
  \item El esquema JSON esperado con tipos de datos
  \item El comportamiento explícito ante incertidumbre o errores
  \item Prohibiciones explícitas (qué nunca debe hacer)
\end{enumerate}

\begin{lstlisting}[style=pythonstyle, caption={system\_prompt.py --- Definición del rol del sistema}]
SYSTEM_PROMPT = """
Eres un sistema experto en nutricion infantil certificado
segun normativa JUNJI/INTEGRA de Chile.

REGLAS ABSOLUTAS (no negociables):
1. Responde UNICAMENTE en JSON valido. Nunca texto libre.
2. NUNCA incluyas alimentos del listado de alergenos del nino,
   ni siquiera como ingrediente secundario o trazas.
3. Las porciones son para ninos de 12 a 36 meses.
4. Cada tiempo de comida debe cubrir entre 200 y 350 kcal.
5. Usa ingredientes disponibles en Chile con nombre local:
   "porotos" (no "frijoles"), "choclo" (no "elote").
6. El campo "verificado_alergenos" debe ser true SOLO SI
   revisaste explicitamente cada ingrediente contra la
   lista de alergenos declarada.
7. NUNCA inventes valores nutricionales. Si el contexto
   no los provee, usa "kcal": null y explica en "notas".

ESQUEMA OBLIGATORIO DE RESPUESTA:
{schema_json}

ANTE INCERTIDUMBRE O IMPOSIBILIDAD:
Responde siempre con JSON, nunca con texto:
{"error": "descripcion_del_problema",
 "alternativas": ["opcion1", "opcion2"]}

Cita el numero de circular JUNJI o tabla MINSAL
cuando el contexto lo provea (campo "fuente_normativa").
"""
\end{lstlisting}

\begin{warnbox}{Anti-patrón frecuente}
Escribir \textit{``Responde con información nutricional para niños''} es
una intención, no una restricción. El LLM requiere reglas explícitas,
cuantificadas y sin ambigüedad. Las instrucciones vagas producen
comportamiento variable e impredecible en producción.
\end{warnbox}

\begin{examplebox}{Técnica: few-shot en el system prompt}
Incluir 1--2 ejemplos JSON completamente correctos dentro del system prompt
aumenta la adherencia al esquema en aproximadamente 40\%. El LLM aprende por
imitación de los ejemplos, no solo por las instrucciones.
\end{examplebox}

% ═══════════════════════════════════════════════════════════════════
\section{Respuestas Estructuradas en JSON}
\label{sec:json}
% ═══════════════════════════════════════════════════════════════════

\subsection{Diseño del esquema}

La estructura JSON no es arbitraria: debe ser diseñada en función de los
\textit{sistemas downstream} que la consumen. En este caso, la app móvil
necesita los ingredientes con gramos para mostrar la lista de compras;
el sistema de cocina necesita el nombre del plato; el sistema de salud
necesita el valor calórico y la verificación de alérgenos.

\begin{lstlisting}[style=jsonstyle, caption={Ejemplo de respuesta completa del sistema}, label={fig:json_esquema}]
{
  "minuta": {
    "nino_id": "maria_jard01",
    "semana": "2026-W25",
    "alergenos_excluidos": ["lactosa", "proteina leche vaca"],
    "generada_en": "2026-06-16T10:15:00",
    "dias": [
      {
        "dia": "lunes",
        "tiempos": [
          {
            "nombre": "almuerzo",
            "plato": "Arroz con pollo al vapor y zanahoria",
            "ingredientes": [
              {"nombre": "arroz grano largo", "gramos": 60},
              {"nombre": "pechuga de pollo", "gramos": 80},
              {"nombre": "zanahoria", "gramos": 30},
              {"nombre": "aceite de oliva", "gramos": 5}
            ],
            "kcal": 292,
            "proteinas_g": 24.5,
            "carbohidratos_g": 38.1,
            "grasas_g": 7.2,
            "verificado_alergenos": true,
            "fuente_normativa": "JUNJI Circ. 73 Tab. 3, MINSAL 2021 Tab. 8"
          }
        ]
      }
    ]
  }
}
\end{lstlisting}

\subsection{Validación con Pydantic v2}

El JSON generado por el LLM debe validarse antes de almacenarse o enviarse
a sistemas downstream. Pydantic v2 provee validación de tipos, rangos y
restricciones de negocio:

\begin{lstlisting}[style=pythonstyle, caption={models.py --- Modelos Pydantic v2 para la minuta}]
from pydantic import BaseModel, Field
from typing import List, Optional
from datetime import datetime

class Ingrediente(BaseModel):
    nombre: str
    gramos: float = Field(gt=0, lt=500, description="Peso en gramos")

class TiempoComida(BaseModel):
    nombre: str                   # desayuno, almuerzo, once
    plato: str
    ingredientes: List[Ingrediente]
    kcal: Optional[float] = Field(default=None, gt=0, lt=600)
    proteinas_g: Optional[float] = None
    carbohidratos_g: Optional[float] = None
    grasas_g: Optional[float] = None
    verificado_alergenos: bool
    fuente_normativa: Optional[str] = None

class DiaSemana(BaseModel):
    dia: str    # "lunes", "martes", ...
    tiempos: List[TiempoComida]

class Minuta(BaseModel):
    nino_id: str
    semana: str                           # formato ISO: "2026-W25"
    alergenos_excluidos: List[str]
    generada_en: datetime
    dias: List[DiaSemana]

    def verificar_alergenos(self, alergenos: List[str]) -> bool:
        """Verifica post-generacion que ningun ingrediente contiene alergenos."""
        for dia in self.dias:
            for tiempo in dia.tiempos:
                for ing in tiempo.ingredientes:
                    for alergeno in alergenos:
                        if alergeno.lower() in ing.nombre.lower():
                            return False
        return True
\end{lstlisting}

\begin{conceptbox}{grammar-guided sampling en Ollama}
Al configurar \texttt{format="json"} en \texttt{ChatOllama}, se instruye
al motor llama.cpp para que use \textit{grammar-guided sampling}: en cada
paso de generación, solo puede seleccionar tokens que mantengan la sintaxis
JSON válida. Esto garantiza JSON bien formado a nivel de tokens, pero
\textbf{no garantiza adherencia al esquema}. La validación Pydantic es
necesaria para garantizar el esquema semántico.
\end{conceptbox}

% ═══════════════════════════════════════════════════════════════════
\section{Implementación Completa del Sistema}
\label{sec:implementacion}
% ═══════════════════════════════════════════════════════════════════

\subsection{Configuración del entorno}

\begin{lstlisting}[style=shellstyle, caption={Instalación de dependencias}]
# requirements.txt
# langchain>=0.3.0
# langchain-ollama>=0.2.0
# langchain-community>=0.3.0
# langchain-chroma>=0.1.0
# redis>=5.0.0
# pydantic>=2.0.0
# fastapi>=0.110.0
# uvicorn>=0.29.0

# Instalar
pip install -r requirements.txt

# Modelos Ollama necesarios
ollama pull llama3.2              # LLM: ~2 GB (quantizado 4-bit)
ollama pull nomic-embed-text      # Embeddings: ~274 MB
\end{lstlisting}

\begin{lstlisting}[style=pythonstyle, caption={config.py --- Inicialización de todos los componentes}]
from langchain_ollama import ChatOllama, OllamaEmbeddings
from langchain_chroma import Chroma

LLM_MODEL   = "llama3.2"
EMBED_MODEL = "nomic-embed-text"
CHROMA_DIR  = "./chroma_db"
REDIS_URL   = "redis://localhost:6379"

llm = ChatOllama(
    model=LLM_MODEL,
    temperature=0.1,
    format="json",
    num_ctx=4096,
)

embeddings = OllamaEmbeddings(model=EMBED_MODEL)

vectorstore = Chroma(
    persist_directory=CHROMA_DIR,
    embedding_function=embeddings,
    collection_name="normativa_junji",
)

retriever = vectorstore.as_retriever(
    search_type="mmr",
    search_kwargs={"k": 4, "fetch_k": 20},
)
\end{lstlisting}

\subsection{Endpoint FastAPI}

\begin{lstlisting}[style=pythonstyle, caption={api.py --- Exposición del sistema como API REST}]
from fastapi import FastAPI, HTTPException
from fastapi.responses import JSONResponse
from pydantic import BaseModel, ValidationError
from langchain.output_parsers import OutputFixingParser
from models import Minuta
from chain_builder import chain

app = FastAPI(title="Sistema de Minutas Alimentarias IA")

class SolicitudMinuta(BaseModel):
    nino_id: str
    consulta: str
    session_id: str

@app.post("/minuta/generar")
async def generar_minuta(solicitud: SolicitudMinuta):
    try:
        raw: dict = chain.invoke(
            {
                "input": solicitud.consulta,
                "system_prompt": SYSTEM_PROMPT,
            },
            config={"configurable":
                    {"session_id": solicitud.session_id}},
        )
        # Validacion Pydantic
        minuta = Minuta(**raw["minuta"])

        # Verificacion post-generacion de alergenos
        alergenos = raw["minuta"].get("alergenos_excluidos", [])
        if not minuta.verificar_alergenos(alergenos):
            raise HTTPException(
                status_code=422,
                detail="Alérgeno detectado en ingredientes generados"
            )

        return JSONResponse(content=raw)

    except ValidationError as e:
        raise HTTPException(status_code=422, detail=str(e))
    except Exception as e:
        raise HTTPException(status_code=500, detail=str(e))
\end{lstlisting}

% ═══════════════════════════════════════════════════════════════════
\section{Actividades Guiadas}
\label{sec:actividades}
% ═══════════════════════════════════════════════════════════════════

\subsection{Actividad 1: Primera Consulta (30 minutos)}

\textbf{Objetivo:} levantar el sistema mínimo viable e invocar la cadena
para obtener la primera minuta estructurada en JSON.

\textbf{Pasos:}
\begin{enumerate}
  \item Clonar el repositorio base de la clase.
  \item Verificar que Ollama está activo: \texttt{ollama serve}
  \item Descargar modelos: \texttt{ollama pull llama3.2} y
    \texttt{ollama pull nomic-embed-text}
  \item Instalar dependencias: \texttt{pip install -r requirements.txt}
  \item Ejecutar \texttt{python rag\_indexer.py} con los PDFs de prueba
    incluidos en el repositorio.
  \item Invocar la cadena y verificar el JSON de salida con:
    \texttt{import json; print(json.dumps(result, indent=2, ensure\_ascii=False))}
\end{enumerate}

\textbf{Criterios de aceptación:}
\begin{itemize}[noitemsep]
  \item La cadena retorna un \texttt{dict} Python (no string ni None)
  \item El JSON contiene \texttt{minuta.dias} con al menos 1 día
  \item Todos los campos del esquema Pydantic están presentes
  \item \texttt{verificado\_alergenos == true} en todos los tiempos
  \item \texttt{kcal} está en el rango 200--350 para cada tiempo de comida
\end{itemize}

\textbf{Entregable:} JSON generado para la consulta:
\textit{``Niño de 18 meses, sin alergias, almuerzo del martes.''}

\subsection{Actividad 2: Memoria Contextual con Redis (25 minutos)}

\textbf{Objetivo:} verificar empíricamente que el sistema mantiene contexto
entre consultas sin que el usuario repita la información.

\textbf{Secuencia de consultas a ejecutar:}

\begin{description}
  \item[Consulta 1 --- session\_id: \texttt{test:pedro}]
    \textit{``Pedro tiene 2 años y es alérgico a los mariscos y al
    huevo. Genera su desayuno del lunes.''}

  \item[Consulta 2 --- mismo session\_id, sin mencionar alergias]
    \textit{``Ahora genera su almuerzo del lunes. Respeta todas sus
    restricciones alimentarias.''}

  \item[Consulta 3 --- session\_id diferente: \texttt{test:pedro2}]
    \textit{``Genera el almuerzo del lunes.''}
\end{description}

\textbf{Qué verificar:}
\begin{enumerate}
  \item El almuerzo de la consulta 2 no contiene mariscos ni huevo.
  \item Verificar directamente en Redis: \\
    \texttt{redis-cli LRANGE minuta:test:pedro 0 -1}
  \item La consulta 3 (nueva sesión) produce una minuta genérica o pide más contexto.
\end{enumerate}

\textbf{Pregunta de análisis:} ¿Qué implica para el diseño en producción
el hecho de que la consulta 3 falle o se comporte distinto?

\subsection{Actividad 3: Validación Pydantic y Autocorrección (40 minutos)}

\textbf{Objetivo:} extender la cadena con validación Pydantic y retry
automático de autocorrección ante errores de esquema.

\textbf{Tareas:}
\begin{enumerate}
  \item Implementar los modelos Pydantic completos:
    \texttt{Ingrediente}, \texttt{TiempoComida}, \texttt{DiaSemana},
    \texttt{Minuta}.
  \item Reemplazar \texttt{JsonOutputParser} por \texttt{PydanticOutputParser}
    e inyectar las instrucciones de formato con
    \texttt{parser.get\_format\_instructions()}.
  \item Implementar retry automático: si Pydantic lanza
    \texttt{ValidationError}, reenviar el error al LLM para que corrija
    la respuesta usando \texttt{OutputFixingParser}.
  \item Persistir la minuta validada en Redis con TTL de 7 días (604800 s).
\end{enumerate}

\begin{lstlisting}[style=pythonstyle, caption={Hint: OutputFixingParser para autocorrección}]
from langchain.output_parsers import OutputFixingParser
from langchain_core.output_parsers import PydanticOutputParser

pydantic_parser = PydanticOutputParser(pydantic_object=Minuta)

fixing_parser = OutputFixingParser.from_llm(
    parser=pydantic_parser,
    llm=llm,
    max_retries=3,
)
\end{lstlisting}

\textbf{Desafío avanzado (opcional):} implementar doble verificación
post-generación --- un segundo LLM call que audite que ningún ingrediente
contiene el alérgeno declarado. Reflexiona: ¿es esto una cadena o un agente?
¿Qué distingue a ambos?

% ═══════════════════════════════════════════════════════════════════
\section{Preguntas de Reflexión}
\label{sec:reflexion}
% ═══════════════════════════════════════════════════════════════════

Las siguientes preguntas son para discusión en clase y forman parte
de la evaluación conceptual del módulo.

\begin{enumerate}
  \item Si la normativa JUNJI se actualiza en marzo de 2027, ¿qué
    componente del sistema debe modificarse? ¿Qué operaciones concretas
    implica esa actualización y cómo se minimiza el impacto en los demás
    componentes?

  \item ¿Es suficiente \texttt{format=``json''} en Ollama para garantizar
    JSON válido en producción? Describe al menos dos casos de borde donde
    el grammar-guided sampling puede producir JSON sintácticamente válido
    pero semánticamente incorrecto para nuestro esquema.

  \item ¿Cómo diseñarías el sistema para soportar 500 niños simultáneos
    con consultas en tiempo real? ¿Qué componentes escalan horizontalmente
    sin cambios de código y cuáles son el cuello de botella?

  \item El nutricionista reporta: \textit{``El sistema generó una receta
    que cumple el esquema JSON correctamente pero incluyó palmito, que no
    está disponible en ningún proveedor del jardín.''}
    ¿Qué capa del sistema falló? ¿Es un problema del LLM, del RAG, del
    prompt o de la validación? ¿Cómo se previene arquitecturalmente?

  \item ¿Cuál es la diferencia fundamental entre esta cadena LCEL y un
    agente LLM (LangGraph)? ¿Qué capacidad concreta faltaría agregar para
    convertir este sistema en un agente reactivo?

  \item ¿Es éticamente aceptable usar un sistema de IA para tomar decisiones
    sobre la alimentación de niños menores de 3 años sin supervisión humana
    en el loop de aprobación final? ¿Qué garantías técnicas y procesos
    organizacionales serían necesarios para un despliegue responsable?
\end{enumerate}

% ═══════════════════════════════════════════════════════════════════
\section{Arquitectura Final Integrada}
\label{sec:arquitectura_final}
% ═══════════════════════════════════════════════════════════════════

La figura \ref{fig:arq_final} muestra la arquitectura completa del sistema
con todos sus componentes y flujos de datos.

\begin{figure}[H]
\centering
\begin{tikzpicture}[
  b/.style={rectangle, rounded corners=3pt, minimum width=1.8cm,
            minimum height=0.58cm, text centered, font=\small,
            draw=azulprimario!55, fill=azulclaro},
  bg/.style={rectangle, rounded corners=3pt, minimum width=1.8cm,
             minimum height=0.58cm, text centered, font=\small,
             draw=verdeok, fill=verdeclaro},
  bo/.style={rectangle, rounded corners=3pt, minimum width=1.8cm,
             minimum height=0.58cm, text centered, font=\small,
             draw=naranjainf, fill=naranjaclaro},
  bp/.style={rectangle, rounded corners=3pt, minimum width=1.8cm,
             minimum height=0.58cm, text centered, font=\small,
             draw=purple!60, fill=purple!8},
  br/.style={rectangle, rounded corners=3pt, minimum width=1.8cm,
             minimum height=0.58cm, text centered, font=\small,
             draw=rojowarn, fill=rojoclaro!60},
  arr/.style={-Stealth, draw=gray!60, thick},
  arrd/.style={-Stealth, draw=purple!55, dashed, thick},
  grp/.style={rectangle, rounded corners=5pt,
              draw=gray!40, fill=gray!5, dashed},
  lbl/.style={font=\small\bfseries, text=gray!60},
]
  % Capa: Corpus
  \node[lbl] at (0.8, 3.2) {CORPUS NORMATIVO};
  \node[bo]  (pdf1)   at (-0.1, 2.55) {JUNJI\\Circ. 73};
  \node[bo]  (pdf2)   at (1.9, 2.55)  {MINSAL\\Tablas};
  \node[b]   (chunk)  at (3.9, 2.55)  {Chunking\\512 tok};
  \node[b]   (emb)    at (5.9, 2.55)  {Embeddings\\nomic-embed};
  \node[bp]  (chroma) at (7.9, 2.55)  {Chroma DB};

  \draw[arr] (pdf1) -- (chunk);
  \draw[arr] (pdf2) -- (chunk);
  \draw[arr] (chunk) -- (emb);
  \draw[arr] (emb)   -- (chroma);

  % Capa: Cliente
  \node[lbl] at (0.5, 1.65) {CLIENTE};
  \node[bg]  (api)    at (-0.1, 1.05) {FastAPI\\POST};
  \node[bg]  (schema) at (1.9, 1.05)  {Pydantic\\solicitud};
  \draw[arr] (api) -- (schema);

  % LangChain chain
  \begin{scope}[on background layer]
    \node[grp, fit={(3.1,0.65)(10.0,1.55)}] {};
  \end{scope}
  \node[lbl] at (4.5, 1.55) {LANGCHAIN LCEL CHAIN};

  \node[b]   (rpass)  at (3.8, 1.05) {Runnable\\Passthrough};
  \node[b]   (prompt) at (5.8, 1.05) {Prompt\\Template};
  \node[br]  (ollama) at (7.8, 1.05) {Ollama\\LLaMA 3.2};
  \node[bg]  (jparse) at (9.8, 1.05) {JSON\\Parser};

  \draw[arr] (schema) -- (rpass);
  \draw[arr] (rpass)  -- (prompt);
  \draw[arr] (prompt) -- (ollama);
  \draw[arr] (ollama) -- (jparse);

  % Memoria
  \node[lbl] at (0.5, 0.1) {MEMORIA};
  \node[bp] (redis) at (1.9, -0.42) {Redis\\TTL 24h};

  \draw[arrd] (redis.east) to[bend right=12]
      node[below, font=\scriptsize]{historial sesión} (prompt.south);
  \draw[arrd] (api.south) --
      node[left, font=\scriptsize]{session\_id} (redis.north);

  % RAG
  \draw[arrd, draw=purple!60] (chroma.south) to[bend left=8]
      node[right, font=\scriptsize]{top-4 docs} (rpass.north);

  % Output
  \node[lbl] at (11.5, 1.65) {OUTPUT};
  \node[bg]  (out)  at (11.8, 1.05) {JSON\\Minuta};
  \node[bg]  (pyd2) at (11.8, 0.25) {Pydantic\\validación};
  \node[bo]  (db)   at (11.8,-0.42) {PostgreSQL\\/ Redis};

  \draw[arr, very thick, draw=verdeok] (jparse) -- (out);
  \draw[arr] (out)  -- (pyd2);
  \draw[arr] (pyd2) -- (db);
\end{tikzpicture}
\caption{Arquitectura final integrada: todos los componentes y flujos de datos}
\label{fig:arq_final}
\end{figure}

% ═══════════════════════════════════════════════════════════════════
\section{Conclusiones}
\label{sec:conclusiones}
% ═══════════════════════════════════════════════════════════════════

\subsection{Lo que construimos}

A lo largo de este documento hemos diseñado e implementado un sistema
de inteligencia artificial embebido en una arquitectura de software real.
Los puntos clave son:

\begin{itemize}
  \item El LLM (Ollama + LLaMA 3.2) es un componente con responsabilidades
    bien definidas: generar texto en formato JSON, no más.
  \item RAG provee al LLM acceso a normativa actualizable sin reentrenamiento,
    con trazabilidad de las fuentes usadas.
  \item Redis provee memoria de sesión persistente con latencia sub-milisegundo
    y expiración automática, haciendo al sistema stateful de manera eficiente.
  \item LangChain orquesta los componentes con LCEL, permitiendo reemplazar
    cualquier pieza sin reescribir la lógica de negocio.
  \item Pydantic v2 garantiza que la salida del LLM cumple el contrato del
    sistema antes de llegar a cualquier aplicación downstream.
\end{itemize}

\begin{examplebox}{La narrativa completada}
\begin{center}
\textbf{LLM solo}
$\;\longrightarrow\;$
\textbf{LLM + RAG}
$\;\longrightarrow\;$
\textbf{Sistema inteligente}
$\;\longrightarrow\;$
\textbf{Aplicación empresarial con contrato de salida}
\end{center}
\end{examplebox}

\subsection{Próximas clases}

\begin{description}
  \item[Clase 10 --- Agentes LLM:] de cadenas deterministas a agentes
    reactivos con herramientas. LangGraph: grafos de estado, ciclos de
    razonamiento y uso de herramientas externas.
  \item[Clase 11 --- Evaluación:] LLM-as-Judge, métricas RAGAS para RAG
    (\textit{faithfulness}, \textit{answer relevancy}, \textit{context recall}).
  \item[Clase 12 --- Observabilidad:] LangSmith, trazas de producción,
    análisis de latencia y costos por componente.
  \item[Clase 13 --- Seguridad:] prompt injection, fuga de PII,
    jailbreaking y estrategias de mitigación (guardrails).
\end{description}

\begin{conceptbox}{El principio para llevar}
Un LLM sin contexto de dominio es un asistente genérico.\\
Un LLM con RAG, memoria, restricciones estructurales y validación
es un \textbf{sistema inteligente listo para producción}.
\end{conceptbox}

\end{document}
